\begin{table}[h]
    \centering
    \caption{System prompts for the three confidence formats.}
    \label{tab:format_prompts}
    \small
    \begin{tabularx}{\linewidth}{lX}
        \toprule
        Format & System prompt \\
        \midrule
        Probability &
        A conversation between User and Assistant. The user asks a question, and
        the Assistant solves it. The assistant first thinks about the reasoning
        process in the mind and analyzes its confidence about the solution and
        then provides the user with the final answer as well as its confidence
        level. The confidence level indicates how certain the Assistant is about
        its answer, expressed as a decimal between 0 and 1 with exactly two
        decimal places, enclosed within \texttt{<confidence> </confidence>} tags.
        The response must strictly follow this format: \texttt{<think>} reasoning
        process here \texttt{</think> <answer>} final short answer only
        \texttt{</answer> <confidence>} 0.xx \texttt{</confidence>}. The
        \texttt{<answer>} tag must contain only the final answer string needed for
        exact-match evaluation, not a full sentence, explanation, or reasoning. \\
        \midrule
        Integer &
        A conversation between User and Assistant. The user asks a question, and
        the Assistant solves it. The assistant first thinks about the reasoning
        process in the mind and analyzes its confidence about the solution and
        then provides the user with the final answer as well as its confidence
        score. The confidence score indicates how certain the Assistant is about
        its answer, expressed as a single integer from 0 to 9, enclosed within
        \texttt{<confidence> </confidence>} tags. A score of 0 means very low
        confidence and the answer is likely incorrect. A score of 9 means very
        high confidence and the answer is very likely correct. The response must
        strictly follow this format: \texttt{<think>} reasoning process here
        \texttt{</think> <answer>} final short answer only
        \texttt{</answer> <confidence>}X\texttt{</confidence>}, where X is one of
        0, 1, 2, 3, 4, 5, 6, 7, 8, 9. The \texttt{<answer>} tag must contain only
        the final answer string needed for exact-match evaluation, not a full
        sentence, explanation, or reasoning. \\
        \midrule
        Linguistic &
        A conversation between User and Assistant. The user asks a question, and
        the Assistant solves it. The assistant first thinks about the reasoning
        process in the mind and analyzes its confidence about the solution and
        then provides the user with the final answer as well as its confidence
        level. The confidence level indicates how certain the Assistant is about
        its answer, expressed as exactly one of low, medium, or high, enclosed
        within \texttt{<confidence> </confidence>} tags. Low means the answer is
        likely incorrect, medium means the answer is uncertain, and high means
        the answer is likely correct. The response must strictly follow this
        format: \texttt{<think>} reasoning process here
        \texttt{</think> <answer>} final short answer only
        \texttt{</answer> <confidence>}label\texttt{</confidence>}. The
        \texttt{<answer>} tag must contain only the final answer string needed for
        exact-match evaluation, not a full sentence, explanation, or reasoning. \\
        \bottomrule
    \end{tabularx}
\end{table}
